\section{Introdução e Objetivo}

O reconhecimento automático de dígitos manuscritos (de 0 a 9) é um problema clássico de visão computacional e aprendizado de máquina, amplamente utilizado para validar modelos de classificação multiclasse. Nesse cenário, o conjunto de dados MNIST consolida-se como referência padrão, reunindo 70.000 imagens em escala de cinza com resolução de $28 \times 28$ pixels. Quando vetorizada, cada imagem corresponde a um vetor de 784 atributos numéricos, permitindo que modelos conexionistas aprendam fronteiras de decisão não lineares diretamente a partir das intensidades de pixel, sem extração manual de características.

O objetivo deste trabalho é implementar do zero, em Python com NumPy (\texttt{learn.py}), uma Rede Neural Artificial do tipo Perceptron Multicamadas (MLP) de duas camadas, composta por uma camada escondida com ativação ReLU e uma camada de saída com função Softmax. Especificamente, busca-se: (i)~implementar a inicialização de parâmetros, as funções de ativação, a propagação direta, a perda de entropia cruzada, a atualização por gradiente descendente e a métrica de acurácia; (ii)~avaliar sistematicamente o impacto da taxa de aprendizado, do número de épocas e da largura da camada escondida sobre a perda de treinamento, a acurácia de teste e o tempo computacional; e (iii)~obter uma configuração com acurácia de teste $\ge 90\%$, validando qualitativamente seu comportamento na interface interativa de desenho (\texttt{server.py}).
